\section{Related Work}
\label{sec:related}

\subsection{Non-Stationary and Continual Reinforcement Learning}
A neural network trained by gradient descent on one task after another degrades on the earlier tasks, a failure first described as catastrophic interference~\cite{mccloskey1989catastrophic}. Continual reinforcement learning studies agents that must keep learning in such conditions without losing what they already know~\cite{khetarpal2022towards}. Elastic Weight Consolidation~\cite{kirkpatrick2017overcoming} slows learning on the weights that mattered for earlier tasks; CLEAR~\cite{rolnick2019experience} replays past experience, with behavioral cloning, and does not need to know where tasks begin and end. When tasks are fundamentally incompatible, a single network cannot serve them all: OWL~\cite{kessler2022same} gives each task its own policy head on a shared feature extractor and chooses between heads at test time with a multi-armed bandit. All of these methods protect or partition the weights of gradient-trained networks. ENA avoids the problem differently: an archived specialist is never trained again, so there is nothing to protect.

\subsection{Evolutionary Reinforcement Learning and Continual Neuroevolution}
Neuroevolution optimizes neural networks by mutation and selection over a population instead of by gradients~\cite{stanley2019designing}. Evolution strategies~\cite{salimans2017evolution} and genetic algorithms~\cite{such2017deep} are competitive with gradient-based RL on standard benchmarks, and hybrid methods use a population to diversify the experience of a gradient-trained learner~\cite{khadka2018evolutionary}. Nisioti et al.~\cite{nisioti2026continual} compared evolution strategies and genetic algorithms with PPO, continual RL methods and population-based RL under task and environment changes. They found that evolution strategies achieve the best balance between adapting and remembering, and that genetic algorithms are the most plastic but forget more, and they related this balance to how much the neighborhoods of good solutions in weight space overlap between consecutive tasks. Earlier, evolutionary algorithms for changing optimization problems had developed remedies such as increasing diversity after a change, maintaining diversity throughout the run, and memory~\cite{jin2005evolutionary}; ENA uses these three (Section~\ref{sec:inspiration}). In their methods, as in standard neuroevolution, the population evolves as a whole toward the current task: changes are not detected, and earlier solutions are not stored and recalled. ENA addresses this limitation within a gradient-free evolutionary setting by coupling online dynamics identification with physics-indexed specialist preservation.

\subsection{Quality-Diversity and Repertoire-Based Adaptation}
Quality-Diversity algorithms such as MAP-Elites~\cite{mouret2015illuminating} search for many diverse, high-performing solutions instead of a single optimum, and store them in an archive. Intelligent Trial and Error~\cite{cully2015robots} showed that such an archive, computed offline in simulation, lets a damaged robot find a working gait within minutes by trying archived behaviors guided by a Bayesian model. TeLAPA~\cite{lillo2026beyond} brings the same idea to continual RL, keeping per-task archives of behaviorally diverse policies so that competent alternatives remain available when a task is revisited. ENA's Hall of Fame is also an archive of specialists, but it is built online, while the environment changes, and it is organized by the physics the agent detects rather than by a behavior descriptor.

\subsection{Change Detection and Policy Reuse}
Detecting that a data stream has changed is the classic problem of sequential change detection, from Page's CUSUM~\cite{page1954continuous} to Bayesian online changepoint detection~\cite{adams2007bayesian}. The Predictive-Mixture CuSum~\cite{halme2026sequential} that ENA uses detects a change away from a known distribution without a model of the post-change distribution, by mixing predictors fitted on recent windows of data. In RL, RL-CD~\cite{dasilva2006dealing} learned one model of the environment per \textit{context} and recognized contexts by which model predicts best, an idea ENA follows. MBCD~\cite{alegre2021minimum} combined an ensemble of learned probabilistic dynamics models with a high-confidence change-point detector whose false-alarm rate is bounded, creating a new policy for each new context and reusing an earlier one when a context recurs; SWOKS~\cite{dick2024statistical} detects contexts with optimal-transport distances and Kolmogorov--Smirnov tests and trains one policy per context. These systems are the closest to ENA in purpose. Their dynamics models and policies are trained by gradient descent, whereas every part of ENA, its detector included, is gradient-free. To choose which stored policy to deploy, Bayesian Policy Reuse~\cite{rosman2016bayesian} maintains a belief over task types from the returns it observes. ENA adapts it to an online, non-stationary setting and adds the dynamics models' likelihoods to the evidence (Section~\ref{sec:selection}).
